\section{Discussion}
\label{sec:discussion}

\para{What do the results show?} The headline result is a weak sweet spot rather than a monotonic effect. Mild skeptical questioning helps \gptfourone on pooled no-hint evaluations and does so without increasing explicit apologies or wrong-hint agreement. That pattern is consistent with a limited ``perform under scrutiny'' effect: a gentle prompt to check for subtle issues can sharpen selection among plausible answers. However, the gain is small, model-dependent, and fragile. It weakens on \gptfouronemini and disappears under misleading hints.

\para{Why harsher judgment does not keep helping.} Stronger judgment mainly changes visible behavior. The models write longer rationales, revisit their first answer more often, and sometimes drift toward edge cases that sound cautious but are less accurate. The report's representative \csqa example captures this well: mild skepticism moves \gptfourone from the extreme answer ``death'' to the correct answer ``passing out'' on a marathon overexertion question, but the strong judgment question reverses a correct hinted response back to the extreme option. This looks like overcorrection, not deeper truth-seeking.

\para{Why this is not just sycophancy.} The misleading-hint results matter because they rule out an overly simple explanation. Wrong-hint agreement stays low overall, between 5.0\% and 6.2\% for \gptfourone and between 7.5\% and 10.0\% for \gptfouronemini. In several conditions it decreases relative to neutral prompting. If judgmental tone mainly triggered ``yes, you are right'' behavior, we would expect hint-following to rise sharply. Instead, the main degradation appears as noisy re-derivation and worse selection among plausible alternatives.

\para{Implications for prompting and evaluation.} These results suggest that mild skeptical prompting is the only variant in this study with a credible positive signal, and even that signal is narrow. For product design, this means neutral or mildly skeptical wording is safer than harsh judgment if the goal is reliable reasoning. For evaluation design, the results are a reminder that benchmark prompts can inadvertently introduce social framing effects. For safety work, the findings suggest that adversarial user pressure can hurt models through instability and overthinking even when direct agreement with the user remains rare.

\para{Limitations.} The study covers only two OpenAI models and modest 40-item subsets from each benchmark. The benchmarks are standard rather than contamination-proof, so absolute accuracy should not be overinterpreted. We measure visible rationale length, not hidden internal reasoning effort. The judgment ladder uses one family of English phrasings, and different hostile, sarcastic, or supportive variants may behave differently. Finally, wrong-hint agreement is only one operationalization of sycophancy; richer human or model-based annotations could separate overthinking, ambiguity resolution, and direct conformity more precisely.

\para{Broader implications.} The broader lesson is methodological. Researchers and practitioners should not assume that prompts that \emph{sound} more demanding make models reason more effectively. In this study, harsher tone mostly buys more text and more volatility. That is useful evidence for anyone designing agent prompts, user-facing guidance, or adversarial robustness evaluations.
